% ============================================================================
% config.tex  --  QUI SI MODIFICANO I DATI DELLA TESI
% ============================================================================
% Compila questo file per primo: frontespizio, metadati del PDF (main.xmpdata)
% e interruttori leggono tutto da qui. Non servono modifiche al frontespizio.
%
% ATTENZIONE: nei valori evita comandi LaTeX complessi (soprattutto in
% \TitoloTesi e \Candidato, che finiscono anche nei metadati del PDF).
% Le lettere accentate e l'apostrofo vanno bene; per le virgolette usa ' e ".
% ============================================================================

% ----------------------------------------------------------------------------
% Dati del frontespizio e dei metadati del PDF
% ----------------------------------------------------------------------------
\newcommand{\TitoloTesi}{Titolo della tesi: sottotitolo se presente}
\newcommand{\TipoTesi}{Tesi sperimentale di specializzazione}
\newcommand{\TipoCorso}{Scuola di Specializzazione in}
\newcommand{\NomeCorso}{Nome della Scuola di Specializzazione}
\newcommand{\Dipartimento}{Dipartimento di Nome del Dipartimento}

\newcommand{\Relatore}{Prof. Nome Cognome}
% Correlatore: lascia vuoto (\newcommand{\Correlatore}{}) se non c'e'.
\newcommand{\Correlatore}{}
\newcommand{\Candidato}{Dott. Nome Cognome}
\newcommand{\Matricola}{XX0000000}
\newcommand{\AnnoAccademico}{AAAA/AAAA}

% Metadati PDF aggiuntivi (usati da main.xmpdata). Le parole chiave si
% separano con \sep, come nell'esempio.
\newcommand{\ParoleChiaveTesi}{Parola chiave 1\sep Parola chiave 2\sep Parola chiave 3}

% Logo dell'Ateneo (vedi FIGURE/README.md). Per rimuoverlo dal frontespizio
% basta lasciare vuoto il valore: \newcommand{\FileLogo}{}
% ATTENZIONE PDF/A: il file deve essere PNG/JPG/PDF senza canale alpha.
\newcommand{\FileLogo}{FIGURE/logo_unicampania.png}

% ----------------------------------------------------------------------------
% Interruttori (\...true attiva, \...false disattiva)
% ----------------------------------------------------------------------------
% PDF/A-1b: formato per l'archiviazione a lungo termine, richiesto da alcune
% biblioteche d'Ateneo. VERIFICA con la tua biblioteca se serve davvero.
% Se lo disattivi, main.xmpdata viene ignorato e si usa un PDF normale.
\newif\ifpdfa
\pdfatrue

% Pagine facoltative. Metti \...false per toglierle dal PDF.
% Il testo si scrive nei rispettivi file.

% Epigrafe: pagina di citazioni dopo il frontespizio (epigrafe.tex).
\newif\ifepigrafe
\epigrafetrue

% Dedica: pagina con una dedica, dopo l'epigrafe (dedica.tex).
\newif\ifdedica
\dedicatrue

% Ringraziamenti: in fondo alla tesi, dopo la bibliografia (ringraziamenti.tex).
\newif\ifringraziamenti
\ringraziamentitrue

% Elenco delle abbreviazioni: dopo l'indice e gli elenchi di figure e tabelle
% (abbreviazioni.tex).
\newif\ifabbreviazioni
\abbreviazionitrue
